% GENERATED FILE. Edit canonical lessons, then run npm run generate:books.
% canonical-source-hash: fnv1a64:fed5c9135f7adb15
% canonical-lessons: UR-C66-suraj, UR-C66-chad, UR-C66-samandar, UR-C66-ag, UR-C66-barf

\chapter{Sun, Moon, Sea, Fire, Snow}
\label{ch:ur-sun-moon-sea-fire-snow}
\hlchaptermodality{66}

\begin{chapteropening}
\textbf{By the end of this chapter:} \emph{I can say the sun, the moon, the sea, fire and snow.}

Complete the last lesson of chapter 66: I can say the sun, the moon, the sea, fire and snow.
\end{chapteropening}

\section[\texorpdfstring{\ur{سورج} (sūraj)}{سورج (sūraj)}]{\ur{سورج} (sūraj) — the sun}

\label{lesson:UR-C66-suraj}



\begin{quote}
Before the new one: say the Urdu for a star, then the Urdu for a cloud.
\end{quote}

\subsection*{You'll want to know: \ur{سورج}}
\textbf{\ur{سورج}} — \emph{sūraj} — \textquotedblleft{}the sun\textquotedblright{}.

\begin{cousinweb}
\textbf{\ur{سورج}} (\emph{sūraj}) is usually traced to Sanskrit \emph{sūrya}, \textquotedblleft{}sun\textquotedblright{}. CDIAL lists that word, but no Hindi or Urdu form under it (CDIAL 13574).
\end{cousinweb}

\begin{grammarlens}[title={What you've built}]
One more word for this chapter.
\end{grammarlens}

\subsection*{Your turn}
\begin{itemize}
\raggedright
  \item \emph{Say it:} \emph{sūraj}
  \item \emph{Say it:} \emph{sūraj}, once more
  \item \emph{Say it:} say \emph{sūraj}
  \item \emph{Recall:} say the Urdu for a star, then the Urdu for a cloud, then say \emph{sūraj} again
\end{itemize}

\begin{tcolorbox}[breakable,colback=teal!4,colframe=teal!35!black,title={Before you move on}]
What does \ur{سورج} mean? (\textquotedblleft{}The sun\textquotedblright{}.) Say it once more.
\end{tcolorbox}
\section[\texorpdfstring{\ur{چاند} (chā̃d)}{چاند (chā̃d)}]{\ur{چاند} (chā̃d) — the moon}

\label{lesson:UR-C66-chad}



\begin{quote}
Before the new one: say the Urdu for a cloud, then the Urdu for the sun.
\end{quote}

\subsection*{You'll want to know: \ur{چاند}}
\textbf{\ur{چاند}} — \emph{chā̃d} — \textquotedblleft{}the moon\textquotedblright{}.

\begin{cousinweb}
\textbf{\ur{چاند}} (\emph{chā̃d}) is a tadbhava, a word inherited from Sanskrit and reshaped by sound change. It comes from \emph{candra}, \textquotedblleft{}moon\textquotedblright{} (CDIAL 4661). CDIAL lists it under the noun \emph{candra}, \textquotedblleft{}moon\textquotedblright{}, built on the adjective \emph{candra}, \textquotedblleft{}shining\textquotedblright{}.
\end{cousinweb}

\begin{grammarlens}[title={What you've built}]
One more word for this chapter.
\end{grammarlens}

\subsection*{Your turn}
\begin{itemize}
\raggedright
  \item \emph{Say it:} \emph{chā̃d}
  \item \emph{Say it:} \emph{chā̃d}, once more
  \item \emph{Say it:} say \emph{chā̃d}
  \item \emph{Recall:} say the Urdu for a cloud, then the Urdu for the sun, then say \emph{chā̃d} again
\end{itemize}

\begin{tcolorbox}[breakable,colback=teal!4,colframe=teal!35!black,title={Before you move on}]
What does \ur{چاند} mean? (\textquotedblleft{}The moon\textquotedblright{}.) Say it once more.
\end{tcolorbox}
\section[\texorpdfstring{\ur{سمندر} (samandar)}{سمندر (samandar)}]{\ur{سمندر} (samandar) — the sea}

\label{lesson:UR-C66-samandar}



\begin{quote}
Before the new one: say the Urdu for the sun, then the Urdu for the moon.
\end{quote}

\subsection*{You'll want to know: \ur{سمندر}}
\textbf{\ur{سمندر}} — \emph{samandar} — \textquotedblleft{}the sea\textquotedblright{}.

\begin{cousinweb}
\textbf{\ur{سمندر}} (\emph{samandar}) goes back to Sanskrit \emph{samudra}, \textquotedblleft{}sea\textquotedblright{}, with an added nasal. CDIAL gives no Hindi or Urdu form in this entry, but it lists Old Awadhi \emph{samuṃda}, \textquotedblleft{}sea\textquotedblright{}, with the same nasal (CDIAL 13236).
\end{cousinweb}

\begin{grammarlens}[title={What you've built}]
One more word for this chapter.
\end{grammarlens}

\subsection*{Your turn}
\begin{itemize}
\raggedright
  \item \emph{Say it:} \emph{samandar}
  \item \emph{Say it:} \emph{samandar}, once more
  \item \emph{Say it:} say \emph{samandar}
  \item \emph{Recall:} say the Urdu for the sun, then the Urdu for the moon, then say \emph{samandar} again
\end{itemize}

\begin{tcolorbox}[breakable,colback=teal!4,colframe=teal!35!black,title={Before you move on}]
What does \ur{سمندر} mean? (\textquotedblleft{}The sea\textquotedblright{}.) Say it once more.
\end{tcolorbox}
\section[\texorpdfstring{\ur{آگ} (āg)}{آگ (āg)}]{\ur{آگ} (āg) — fire}

\label{lesson:UR-C66-ag}



\begin{quote}
Before the new one: say the Urdu for the moon, then the Urdu for the sea.
\end{quote}

\subsection*{You'll want to know: \ur{آگ}}
\textbf{\ur{آگ}} — \emph{āg} — \textquotedblleft{}fire\textquotedblright{}.

\begin{cousinweb}
\textbf{\ur{آگ}} (\emph{āg}) is a tadbhava, a word inherited from Sanskrit and reshaped by sound change. It comes from \emph{agni}, \textquotedblleft{}fire\textquotedblright{} (CDIAL 55).
\end{cousinweb}

\begin{grammarlens}[title={What you've built}]
One more word for this chapter.
\end{grammarlens}

\subsection*{Your turn}
\begin{itemize}
\raggedright
  \item \emph{Say it:} \emph{āg}
  \item \emph{Say it:} \emph{āg}, once more
  \item \emph{Say it:} say \emph{āg}
  \item \emph{Recall:} say the Urdu for the moon, then the Urdu for the sea, then say \emph{āg} again
\end{itemize}

\begin{tcolorbox}[breakable,colback=teal!4,colframe=teal!35!black,title={Before you move on}]
What does \ur{آگ} mean? (\textquotedblleft{}Fire\textquotedblright{}.) Say it once more.
\end{tcolorbox}
\section[\texorpdfstring{\ur{برف} (barf)}{برف (barf)}]{\ur{برف} (barf) — snow; ice}

\label{lesson:UR-C66-barf}



\begin{quote}
Before the new one: say the Urdu for the sea, then the Urdu for fire.
\end{quote}

\subsection*{You'll want to know: \ur{برف}}
\textbf{\ur{برف}} — \emph{barf} — \textquotedblleft{}snow; ice\textquotedblright{}.

\begin{cousinweb}
\textbf{\ur{برف}} (\emph{barf}) is a loan, usually traced to Persian \emph{barf}, \textquotedblleft{}snow, ice\textquotedblright{}.
\end{cousinweb}

\begin{grammarlens}[title={What you've built}]
One more word for this chapter.
\end{grammarlens}

\subsection*{Your turn}
\begin{itemize}
\raggedright
  \item \emph{Say it:} \emph{barf}
  \item \emph{Say it:} \emph{barf}, once more
  \item \emph{Say it:} say \emph{barf}
  \item \emph{Recall:} say the Urdu for the sea, then the Urdu for fire, then say \emph{barf} again
\end{itemize}

\begin{tcolorbox}[breakable,colback=teal!4,colframe=teal!35!black,title={Before you move on}]
What does \ur{برف} mean? (\textquotedblleft{}Snow; ice\textquotedblright{}.) Say it once more.
\end{tcolorbox}
